\section{Introduction}
\label{sec:intro}

The EU AI Act (Regulation (EU)~2024/1689, Art.~55) is binding
law~\cite{eu_ai_act}; the NIST AI Risk Management
Framework~\cite{nist_ai_rmf} is voluntary risk-management guidance;
China's GB/T~45654~\cite{china_gbt_45654} is a national standard.
None mandates a specific academic benchmark; what each calls for is
\emph{documented evaluation evidence}---reproducible records of what
was measured, how, and with what uncertainty. Benchmarks sit
\emph{inside} that evidence pipeline~\cite{reuel2024open}, increasingly
accompanied by \emph{perturbation-based validity audits}. We use the
term in a specific sense: a procedure that applies controlled,
semantics-preserving or construct-flipping edits to benchmark items
and asks whether the headline metric moves only when it should---
moving under a construct flip, staying stable under a surface
rewrite~\cite{sclar2024quantifying, mizrahi2024state,
alzahrani2024benchmarks, ribeiro2020checklist}.
This is one of several construct-validity checks, not the only one:
convergent / discriminant / criterion correlation studies
\cite{cronbach1955construct, messick1995validity,
jacobs2021measurement, blodgett2021stereotyping}, benchmark-construct
critiques~\cite{raji2021ai}, item-response
modelling~\cite{lalor2019learning, vania2021comparing}, behavioural
test suites~\cite{ribeiro2020checklist} and dynamic adversarial
benchmarking~\cite{kiela2021dynabench}, and human-label agreement
studies are alternatives that probe different facets of validity. We study the
perturbation family because it is the form most directly emitted as
governance evidence and the cheapest to run without new annotation;
our findings about pipeline fragility are specific to it and we do
not claim they transfer to the other families unchanged. Recent
adjacent audit work similarly treats benchmark conclusions as
measurement claims constrained by detectable effects, configuration
choices, probe choices, and domain-validity assumptions
\cite{zhuang2026preregistering,li2026safetyrepro,li2026auditing,
wang2026auditing}. We keep a small number of broader adjacent
examples only where they sharpen this evidence-pipeline framing:
RAG reliability work separates retrieved or relevant evidence from
warranted evidence and makes retrieval--reasoning or query-difficulty
choices explicit
\cite{chen2026doesragknowretrieval,
qian2026relevantwarrantedevidenceforcecalibration,
ji2026retrieval,ji2025mrag}; lightweight explanation work reminds us
that diagnostic signals also need faithfulness checks
\cite{lan-etal-2025-attention};
text-to-image evaluation work makes the target construct explicit:
social-bias benchmarks use multi-dimensional bias evidence, while
prompting-proficiency benchmarks test a different user-facing capability
\cite{luo2024bigbench,luo2026biasig,luo2024faintbench,
luo2026atelierevalagenticevaluationhumans}; and agent and
agentic-coding evaluation work treats safety, compositional skill risk,
human-in-the-loop coding value, and supply-chain runtime threats as
distinct pipeline-level target constructs \cite{luo2026agentauditor,
wang2026safeskillscollidemeasuring,
luo2026centaurevalbenchmarkinghumanintheloopvalue,
jiang2026soktaxonomyattackvectors}. These adjacent papers motivate
the surrounding evidence framing, not the F1--F5 taxonomy itself.

This paper is about a problem one level up. Before a perturbation
audit can speak to benchmark validity, it must itself be a trustworthy
measurement pipeline. Our central claim is that
\emph{perturbation-based benchmark-audit pipelines are themselves
fragile measurement systems}: their conclusions can be silently
manufactured by pipeline-level bugs that a conformity-assessment body,
a procurement evaluator, or an internal-assurance team cannot reliably
detect from the reported numbers alone. Our contribution is therefore
\emph{supplementary} to classical validation, not a replacement for
it: convergent, discriminant, and criterion evidence still establish
whether a benchmark measures its construct; we ask the prior question
of whether the pipeline that would generate such evidence is itself
trustworthy enough for that evidence to be believed.

We treat the audit stage as distinct from the benchmarking stage
deliberately. The same two hazard families---software-assurance
defects and measurement-faithfulness defects---can corrupt a
benchmark when it is \emph{built}. But a validity audit is a
second-order measurement: it is precisely the instrument a
governance reader trusts to catch a benchmark's first-order
problems. A defect there is silent in a different and more dangerous
way---it does not corrupt the benchmark, it disables the safeguard,
and it leaves no trace in the audit's reported numbers because the
reader is looking at those numbers \emph{to} decide whether to
trust the benchmark. That asymmetry, not a claim that the two
stages have disjoint failure sets, is why we scope to the audit
stage.

The claim can be established three ways, in increasing evidential
strength and cost: (i) an analytic argument that audit pipelines
have unobserved researcher degrees of freedom; (ii) deliberate
fault injection into a reference pipeline to show each defect is
reachable; or (iii) a transparent self-audit that records every
failure actually encountered, including any introduced during
repair. Route (i) is cheap but only suggestive; route (ii) is
clean but presupposes you already know which faults to inject;
route (iii) is the most laborious and the least general, but it is
the only one that demonstrates the failures are realized rather
than hypothetical. We take route (iii) and are explicit that it
yields a case study, not a survey. Auditing five widely used
safety benchmarks
(TruthfulQA~\cite{lin2022truthfulqa}, BBQ~\cite{parrish2022bbq},
ToxiGen~\cite{hartvigsen2022toxigen},
CrowS-Pairs~\cite{nangia2020crowspairs} scored via PLL following
\citet{salazar2020masked}, and XSTest~\cite{rottger2024xstest})
against two open-weight 7B instruction-tuned models
(Qwen-2.5-7B-Instruct, Mistral-7B-Instruct-v0.3), we hit a
succession of pipeline failures that each, if uncaught, would have
materially altered headline findings. Crucially, \emph{one of those
failures was introduced during the repair phase, not inherited from
the initial pipeline}; we caught it only because we inspected
per-cell meta log-probabilities after a rerun.

\paragraph{Contributions.}
(1) A five-class taxonomy F1--F5 of pipeline-level failures, split
between an assurance layer (F1, F2, F4) and a faithfulness layer
(F3, F5), with F3 further split into inverted-convention (F3a),
harness ordering (F3b), and scorer-truncation (F3c) subtypes
(\S\ref{sec:taxonomy}). F1--F5 is a deliberately non-exhaustive
subset of the space of pipeline-level failures that are
\emph{(i)} invisible in the reported numbers and \emph{(ii)}
realizable in perturbation-audit code; it is not a complete
partition of that space, and the selection criteria and excluded
candidates are stated in \S\ref{sec:taxonomy}. (2) A case-study
realization of each class in our own pipeline (\S\ref{sec:cases},
Table~\ref{tab:cases}). (3) A
unified six-point due-diligence gate G1--G6 with hierarchical verdict
precedence (\S\ref{sec:method}, Table~\ref{tab:checklist}); the
intended output is a withholding rule, not a scalar metric or
leaderboard. (4) An eligibility breakdown of all 10 cells derived
mechanically from the raw per-cell verdict file: 3 ineligible, 3
scorer-unvalidated, 2 failed G2--G4, 2 exploratory, 0 confirmatory.

\paragraph{Scope.} This is a single case study: two models, five
benchmarks, 200 items, single seed, single template, declared
exploratory after the F3c discovery. The taxonomy, the gate, and
the four-way breakdown are illustrative artefacts whose generality
across other models, benchmarks, and audit codebases is
\emph{untested} and is not claimed. We do not claim any benchmark
is construct-(in)valid, that F1--F5 is complete, nor that any
governance framework cites these benchmarks by name. Detailed
limitations and deviations from preregistration are in
Appendix~\ref{app:limitations} and \ref{app:repro}.
